\documentclass[11pt,a4paper]{article}

%% PACHETE MATEMATICE
\usepackage{amsmath,amssymb,amsfonts}
\usepackage{amsthm}
\usepackage{mathpazo}
\usepackage{eulervm}
\usepackage{enumerate}
\usepackage{color}
\usepackage{graphicx}
\usepackage[all]{xy}
\usepackage{xcolor}
\usepackage{framed}
\usepackage[unicode]{hyperref}
\setcounter{MaxMatrixCols}{20}
\usepackage{multicol}
% \usepackage[romanian]{babel}


%% ASPECT
\linespread{1.05}
\usepackage[T1]{fontenc}
\usepackage[utf8x]{inputenc}
\usepackage[margin=2cm]{geometry}
% \usepackage[color]{showkeys}
\usepackage{anyfontsize}
\usepackage{titlesec}
\titleformat*{\section}{\large\bfseries}

% \usepackage{fancyhdr}
% \pagestyle{fancy}
% \rhead{\small \color{gray} Notițe de Adrian Manea}
% \lhead{\small \color{gray} IntProb-FIA, 2017---2018, L. Lemnete \& M. Olteanu}


\usepackage[autostyle = false, style = german]{csquotes}
\usepackage[romanian]{babel}
\newcommand\qq{\enquote}

%% COMENZILE MELE
\renewcommand*{\proofname}{Dem.:}
\newcommand\bb{\mathbb}
\newcommand\dr{\mathrm}
\newcommand\kal{\mathcal}
\newcommand\fk{\mathfrak}
\newcommand\op{\oplus}
\newcommand\ot{\otimes}
\newcommand\ds{\displaystyle}
\newcommand\id{\indent}
\newcommand\nid{\noindent}
\newcommand\seq{\subseteq}
\newcommand\sneq{\subsetneq}
\newcommand\speq{\supseteq}
\newcommand\spneq{\supsetneq}
\newcommand\rar{\rightarrow}
\newcommand\Rar{\Rightarrow}
\newcommand\xrar{\xrightarrow}
\newcommand\lar{\leftarrow}
\newcommand\Lar{\Leftarrow}
\newcommand\xlar{\xleftarrow}
\newcommand\lrar{\leftrightarrow}
\newcommand\Lrar{\Leftrightarrow}
\newcommand\ideal{\trianglelefteq}
\newcommand\ai{\text{ a.\^{i}. }}
\newcommand\sk{(\textit{Schi\c{t}\u{a}}) }
\newcommand\rhup{\rightharpoonup}
\newcommand\rhdn{\rightharpoondown}
\newcommand\lhup{\leftharpoonup}
\newcommand\lhdn{\leftharpoondown}
\newcommand\vid{\emptyset}
\newcommand\wh{\widehat}
\newcommand\ol{\overline}
\newcommand\NN{\mathbb{N}}
\newcommand\RR{\mathbb{R}}
\newcommand\ZZ{\mathbb{Z}}
\newcommand\QQ{\mathbb{Q}}
\newcommand\CC{\mathbb{C}}

%% STILURILE MELE
\newtheoremstyle{dotless}{}{}{\itshape}{}{\bfseries}{:}{ }{}

\theoremstyle{dotless}
\newtheorem{theorem}{Teorem\u{a}}[section]
\newtheorem{proposition}{Propozi\c{t}ie}[section]
\newtheorem{lemma}{Lem\u{a}}[section]
\newtheorem{corollary}{Corolar}[section]
\newtheorem{exercise}{Exerci\c{t}iu}

\newtheoremstyle{dotlessdr}{}{}{}{}{\bfseries}{:}{ }{}
\theoremstyle{dotlessdr}
\newtheorem{example}{Exemplu}[section]
\newtheorem{definition}{Defini\c{t}ie}[section]
\newtheorem{remark}{Observa\c{t}ie}[section]




\begin{document}

\begin{center}
  {\Large\textbf{Parțial calcul integral \\
      917B}}
\end{center}


\textbf{NUMĂRUL 1}

1. Folosind derivarea sub semnul integral, calculați:
\[
  I(a) = \int_0^{\frac{\pi}{2}} \ln(\cos^2 x + m^2 \sin^2 x) dx, m > 0.
\]

\vspace{1cm}

2. Calculați $ \ds\int_0^1 \sqrt{\ln \frac{1}{x}} dx $.

\vspace{1cm}

3. Fie $ \Gamma = \{ (x, y) \in \RR^2 \mid x^2 + y^2 = 4 \} $.

Calculați integrala curbilinie:
\[
  \int_\Gamma \frac{-y}{x^2 + y^2} dx + \frac{x}{x^2 + y^2} dy
\]
\begin{enumerate}[(a)]
\item direct (folosind parametrizarea curbei $ \Gamma $);
\item folosind formula Green-Riemann, pe $ K = \mathrm{Int}(\Gamma) - \{ (0, 0) \} $.
\end{enumerate}

\vspace{1cm}

(*) Calculați volumul corpului mărginit de suprafețele de ecuații:
\[
  \begin{cases}
    z &= x^2 + y^2 - 1 \\
    z & = 4 - x^2 - y^2
  \end{cases}.
\]

\newpage

\textbf{NUMĂRUL 2}

1. Folosind derivarea sub integrală, calculați:
\[
  I(a) = \int_0^1 \frac{\arctan ax}{x \sqrt{1 - x^2}} dx, a > 0.
\]

\vspace{1cm}

2. Calculați integrala:
\[
  \int_1^{e} \frac{1}{x} \ln^3 x \cdot (1 - \ln x)^4 dx.
\]

\vspace{1cm}

3. Calculați integrala:
\[
  I = \int_C y^2 dx + x^2 dy,
\]
unde curba $ C $ este intersecția dintre:
\[
  \begin{cases}
    x^2 + y^2 &= 1 \\
    y &\geq 0
  \end{cases} \quad \text{ și } \quad
  \begin{cases}
    x &= t \in [-1, 1] \\
    y &= 0
  \end{cases}
\]
\begin{enumerate}[(a)]
\item folosind parametrizarea curbei (calcul direct);
\item folosind formula Green-Riemann.
\end{enumerate}

\vspace{1cm}

(*) Calculați volumul corpului delimitat de suprafețele:
\[
  \begin{cases}
    z &= x^2 + y^2 - 1 \\
    z &= 2 - x^2 - y^2
  \end{cases}.
\]



\end{document}
